\section{Vectores}

\subsection{Vectores en el plano y en el espacio tridimensional}

Dos puntos $a$ y $b$ que pertenecen a una recta $R$ determinan un segmento denotado por $\overline{ab}$, como se muestra en la siguiente figura:
\begin{figure}[H]
    \centering
    \includegraphics[width=0.4\linewidth]{Imagenes/U4/fig1.png}
    \caption{}
    \label{fig:4.1}
\end{figure}

Si esos dos puntos se dan con un cierto orden, queda determinado un \Emph{vector}. El primer punto dado, $a$, se llama \Emph{origen} del vector o \Emph{punto de aplicación}, y el segundo, $b$, se denomina \Emph{extremo}. Se llama \Emph{vector geométrico} a todo \emph{segmento orientado}. Además, al vector geométrico de origen $a$ y extremo $b$, lo podemos escribir como $\vec{ab}$, aunque otros autores los denotan simplemente como $\mathbf{ab}$ (en negrita). Geométricamente se lo representa de la siguiente manera:
\begin{figure}[H]
    \centering
    \includegraphics[width=0.4\linewidth]{Imagenes/U4/fig2.png}
    \caption{}
    \label{fig:4.2}
\end{figure}
%El vector puede escribirse en función de sus \Emph{componentes}, como veremos más adelante, haciendo uso de pares ordenados (o ternas ordenadas en el caso del espacio tridimensional). Por ejemplo, un vector $\vv$ que tenga 2 componentes, $x$ e $y$, se puede escribir como $\vv = (x,y)$.

%%%%%%%%%%%%%%%%%%%%%%%%%%%%%%%%%%%%%%%%%%%%%%%%%%%%%%%%%%%%
%%%%%%%%%%%%%%%%%%%%%%%%%%%%%%%%%%%%%%%%%%%%%%%%%%%%%%%%%%%%
%%%%%%%%%%%%%%%%%%%%%%%%%%%%%%%%%%%%%%%%%%%%%%%%%%%%%%%%%%%%


\subsubsection*{Características}

Un vector queda siempre determinado por su:
\begin{itemize}
    \item \Emph{Dirección:} La dirección del vector está determinada por la recta $R$ que lo contiene (llamada recta soporte) o la dirección de sus paralelas.

    \item \Emph{Sentido:} El sentido del vector $\vec{ab}$ está determinado por el orden de los puntos $a$ y $b$, origen y extremo. Y lo indicamos mediante una punta de flecha situada en un extremo, que muestra hacia qué lado de la \Emph{línea de acción} (recta soporte) se dirige el vector.

    \item \Emph{Módulo:} El módulo (o también norma) de un vector es la \Emph{longitud} del vector. Para conocer el módulo de un vector es preciso conocer el origen y el extremo del mismo, porque debemos medir desde su origen hasta su extremo. En el vector $\vec{ab}$ el módulo es la longitud del segmento $\overline{ab}$. Al módulo del vector $\vec{ab}$ lo simbolizamos: $|\vec{ab}|$.
\end{itemize}
Se puede ver en la siguiente figura las características recién descritas:
\begin{figure}[H]
    \centering
    \includegraphics[width=0.4\linewidth]{Imagenes/U4/fig3.png}
    \caption{}
    \label{fig:4.3}
\end{figure}

%%%%%%%%%%%%%%%%%%%%%%%%%%%%%%%%%%%%%%%%%%%%%%%%%%%%%%%%%%%%
%%%%%%%%%%%%%%%%%%%%%%%%%%%%%%%%%%%%%%%%%%%%%%%%%%%%%%%%%%%%
%%%%%%%%%%%%%%%%%%%%%%%%%%%%%%%%%%%%%%%%%%%%%%%%%%%%%%%%%%%%


\subsubsection*{Módulo de un vector}

A partir del conocimiento del módulo, podemos definir:
\begin{itemize}
    \item \Emph{Vector Nulo o Cero:} es aquel vector cuyo módulo es cero y \Emph{no} tiene dirección. Lo denotamos como $\vec{0}$ ($\mathbf{0}$). Simbólicamente se escribe $|\vec{0}| = 0$.

    \item \Emph{Versor o Vector Unitario:} es todo vector $\vec{v_0}$ cuyo módulo es exactamente igual a uno. Es decir: $|\vec{v_0}| = 1$.
\end{itemize}


%%%%%%%%%%%%%%%%%%%%%%%%%%%%%%%%%%%%%%%%%%%%%%%%%%%%%%%%%%%%
%%%%%%%%%%%%%%%%%%%%%%%%%%%%%%%%%%%%%%%%%%%%%%%%%%%%%%%%%%%%
%%%%%%%%%%%%%%%%%%%%%%%%%%%%%%%%%%%%%%%%%%%%%%%%%%%%%%%%%%%%


\subsubsection*{Componentes de un vector}

Introduciendo un \Emph{sistema de coordenadas cartesianas rectangulares}, podemos trabajar con vectores en el plano, que lo denotaremos  $\RR^2=\{(x,y)\ /\ x,y \in\RR\}$ (dos dimensiones) o bien, en el espacio, es decir $\RR^3=\{(x,y,z)\ /\ x,y,z\in\RR\}$ (tres dimensiones). Para hacer referencia simultánea de ambos conjuntos utilizaremos la notación $\VV$.

%%%%%%%%%%%%%%%%%%%%%%%%%%%%%%%%%%%%%%%%%%%%%%%%%%%%%%%%%%%%

\subsubsection*{En el plano $\RR^2$}

Para describir un vector analíticamente en el plano, consideramos un vector con  origen en el origen de coordenadas $\bigO=(0,0)$ y como extremo un punto $P(x_1,y_1)$, es decir $\vv = \vec{OP} =(x_1,y_1)$. Los números $x_1$, $y_1$ se denominan las \Emph{componentes del vector} $\vv$. En general, cualquier vector en el espacio bidimensional puede identificarse con las coordenadas de su extremo. Gráficamente:
\begin{figure}[H]
    \centering
    \includegraphics[width=0.35\linewidth]{Imagenes/U4/fig4.png}
    \caption{}
    \label{fig:4.4}
\end{figure}
En el siguiente gráfico, podemos ver que el vector $\vec{\bigO a}$ tiene componentes 4 y 3, respectivamente, por lo que su notación en función de sus componentes es: $\vec{\bigO a} = (4,3)$.
\begin{figure}[H]
    \centering
    \includegraphics[width=0.35\linewidth]{Imagenes/U4/fig5.png}
    \caption{}
    \label{fig:4.5}
\end{figure}

%%%%%%%%%%%%%%%%%%%%%%%%%%%%%%%%%%%%%%%%%%%%%%%%%%%%%%%%%%%%

\subsubsection*{En el espacio $\RR^3$}

Análogamente, un vector en $\RR^3$, es decir, el espacio tridimensional, puede expresarse en función de 3 componentes ortogonales entre sí $x_1,y_1,z_1$ (\figura{4.6}). Entonces, resulta que $\vv = \vec{\bigO P} = (x_1, y_1, z_1)$. Un ejemplo numérico (\figura{4.7}) es el vector $\vec{a} = \vec{\bigO P} = (3,3,2)$.

\noindent
\begin{minipage}[t]{0.49\textwidth}
    \begin{figure}[H]
        \centering
        \includegraphics[width=0.6\linewidth]{Imagenes/U4/fig6.png}
        \caption{}
        \label{fig:4.6}
    \end{figure}
\end{minipage}
\hfill
\begin{minipage}[t]{0.49\textwidth}
    \begin{figure}[H]
        \centering
        \includegraphics[width=0.6\linewidth]{Imagenes/U4/fig7.png}
        \caption{}
        \label{fig:4.7}
    \end{figure}
\end{minipage}
\vspace{0.2cm}


%%%%%%%%%%%%%%%%%%%%%%%%%%%%%%%%%%%%%%%%%%%%%%%%%%%%%%%%%%%%
%%%%%%%%%%%%%%%%%%%%%%%%%%%%%%%%%%%%%%%%%%%%%%%%%%%%%%%%%%%%
%%%%%%%%%%%%%%%%%%%%%%%%%%%%%%%%%%%%%%%%%%%%%%%%%%%%%%%%%%%%

\subsubsection*{Módulo o norma de un vector en función de sus componentes}
% rever una introduccion a la escritura del vector en funcion de sus componentes

A partir de lo expuesto hasta el momento, veremos cómo calcular el módulo de un vector conociendo sus componentes:

\begin{itemize}
    \item \Emph{En $\RR^2$:}
    
    \begin{minipage}[c]{0.55\textwidth}
        \vspace{-1cm}
        Dado el vector $\vv = \vec{op}$, $|\vv|$ es la hipotenusa del triángulo rectángulo, donde $x_1$ e $y_1$ son los catetos. Por el Teorema de Pitágoras, podemos expresar:
            \begin{equation*}
                |\vv|^2 = x_1^2 + y_1^2 \iff |\vv| = +\sqrt{x_1^2 + y_1^2}
            \end{equation*}
    \end{minipage}
    \hfill
    \begin{minipage}[c]{0.4\textwidth}
        \vspace{-1cm}
        \begin{figure}[H]
            \centering
            \includegraphics[width=0.6\linewidth]{Imagenes/U4/fig8.png}
            \caption{}
            \label{fig:4.8}
        \end{figure}
    \end{minipage}

    \item \Emph{En $\RR^3$:}

    
    \begin{minipage}[c]{0.4\textwidth}
        \begin{figure}[H]
            \centering
            \includegraphics[width=0.8\linewidth]{Imagenes/U4/fig9.png}
            \caption{}
            \label{fig:4.9}
        \end{figure}
    \end{minipage}
    \hfill
    \begin{minipage}[c]{0.55\textwidth}
        Por el Teorema de Pitágoras, la diagonal $\overline{\bigO Q}$ verifica que: $\overline{\bigO Q}^2 = x_1^2 + y_1^2$.

        Luego, para hallar el módulo de $\vv$, aplicamos nuevamente Pitágoras, entonces:
            \begin{align*}
                |\vv|^2 = \overline{\bigO Q}^2 + \overline{QP}^2 &\iff |\vv|^2 = x_1^2 + y_1^2 + z_1^2 \iff \\
                &\iff|\vv| = +\sqrt{x_1^2 + y_1^2 + z_1^2}
            \end{align*}
    \end{minipage}
\end{itemize}


%%%%%%%%%%%%%%%%%%%%%%%%%%%%%%%%%%%%%%%%%%%%%%%%%%%%%%%%%%%%
%%%%%%%%%%%%%%%%%%%%%%%%%%%%%%%%%%%%%%%%%%%%%%%%%%%%%%%%%%%%
%%%%%%%%%%%%%%%%%%%%%%%%%%%%%%%%%%%%%%%%%%%%%%%%%%%%%%%%%%%%


\subsubsection*{Igualdad de vectores}

\Emph{Vector libre:} es el vector que puede trasladarse sin modificar su módulo, dirección y sentido.

\begin{definitionbox}{Igualdad de Vectores}{4.1}
    Dos vectores libres son iguales si tienen igual dirección, sentido y módulo. Los vectores iguales se llaman también \Emph{vectores equipolentes}.
\end{definitionbox}
En el siguiente gráfico se muestra que los vectores $\vv,\ \vec{w},$ y $\vec{z}$ (\figura{4.10}) son equipolentes entre sí, y $\vec{a}$ es equipolente a $\vec{b}$ (\figura{4.11}).

\noindent
\begin{minipage}[b]{0.45\textwidth}
    \begin{figure}[H]
        \centering
        \includegraphics[width=0.3\linewidth]{Imagenes/U4/fig10.png}
        \caption{}
        \label{fig:4.10}
    \end{figure}
\end{minipage}
\hfill
\begin{minipage}[b]{0.45\textwidth}
    \begin{figure}[H]
        \centering
        \includegraphics[width=0.3\linewidth]{Imagenes/U4/fig11.png}
        \caption{}
        \label{fig:4.11}
    \end{figure}
\end{minipage}

\subsubsection*{Igualdad de vectores en función de sus componentes}

Dos vectores de $\VV$ son iguales si se verifica la igualdad de sus respectivas componentes.
\begin{itemize}
    \item \Emph{En $\RR^2$:} Sean $\vu = (x_1,y_1),\ \vv = (x_2,y_2)$ dos vectores en el plano. Resulta que:
        \begin{equation*}
            \vu = \vv \iff x_1 = x_2\ \land\ y_1 = y_2
        \end{equation*}

    \item \Emph{En $\RR^3$:} Sean $\vu = (x_1,y_1,z_1),\ \vv = (x_2,y_2,z_2)$ dos vectores en el plano. Resulta que:
        \begin{equation*}
            \vu = \vv \iff x_1 = x_2\ \land\ y_1 = y_2\ \land\ z_1 = z_2
        \end{equation*}
\end{itemize}

%%%%%%%%%%%%%%%%%%%%%%%%%%%%%%%%%%%%%%%%%%%%%%%%%%%%%%%%%%%%

\subsubsection*{Vectores opuestos}

Un vector $\vv$ y otro $-\vv$ se llaman vectores opuestos si tienen igual dirección y módulo, pero tienen sentidos contrarios. Por lo tanto, los vectores opuestos \Emph{no} son equipolentes.

Si $\vv = (x_1, y_1)\in\RR^2$, se define su opuesto como $-\vv = (-x_1, -y_1)$. De manera análoga se lo puede definir en $\RR^3$.


%%%%%%%%%%%%%%%%%%%%%%%%%%%%%%%%%%%%%%%%%%%%%%%%%%%%%%%%%%%%
%%%%%%%%%%%%%%%%%%%%%%%%%%%%%%%%%%%%%%%%%%%%%%%%%%%%%%%%%%%%
%%%%%%%%%%%%%%%%%%%%%%%%%%%%%%%%%%%%%%%%%%%%%%%%%%%%%%%%%%%%

\begin{preguntabox}[Ejercicio 1:]
    Considerar el vector $\vv = (-4,2)$ y hallar las componentes de otros tres vectores que tengan:
        \begin{enumerate}[label=(\alph*)]
            \item La misma dirección que $\vv$ y el mismo sentido.

            \item La misma dirección que $\vv$ y el sentido opuesto.

            \item El mismo módulo que $\vv$.
        \end{enumerate}
\end{preguntabox}


\begin{preguntabox}[Ejercicio 2:]
    Considerar la igualdad: $|\vec{a} + \vec{b}| = |\vec{a}| + |\vec{b}|$. Encontrar dos vectores que verifiquen la igualdad y dos vectores que no.
\end{preguntabox}
%%%%%%%%%%%%%%%%%%%%%%%%%%%%%%%%%%%%%%%%%%%%%%%%%%%%%%%%%%%%
%%%%%%%%%%%%%%%%%%%%%%%%%%%%%%%%%%%%%%%%%%%%%%%%%%%%%%%%%%%%
%%%%%%%%%%%%%%%%%%%%%%%%%%%%%%%%%%%%%%%%%%%%%%%%%%%%%%%%%%%%


\subsection{Versores fundamentales en $\RR^2$ y $\RR^3$}

% pendiente


%%%%%%%%%%%%%%%%%%%%%%%%%%%%%%%%%%%%%%%%%%%%%%%%%%%%%%%%%%%%
%%%%%%%%%%%%%%%%%%%%%%%%%%%%%%%%%%%%%%%%%%%%%%%%%%%%%%%%%%%%
%%%%%%%%%%%%%%%%%%%%%%%%%%%%%%%%%%%%%%%%%%%%%%%%%%%%%%%%%%%%


\subsection{Álgebra vectorial}

Los vectores operan de distintas maneras entre sí y con elementos de otros conjuntos. En este curso estudiaremos:
\begin{itemize}
    \item La suma entre vectores.

    \item La diferencia entre vectores.

    \item El producto de un vector por un escalar.
\end{itemize}


%%%%%%%%%%%%%%%%%%%%%%%%%%%%%%%%%%%%%%%%%%%%%%%%%%%%%%%%%%%%
%%%%%%%%%%%%%%%%%%%%%%%%%%%%%%%%%%%%%%%%%%%%%%%%%%%%%%%%%%%%
%%%%%%%%%%%%%%%%%%%%%%%%%%%%%%%%%%%%%%%%%%%%%%%%%%%%%%%%%%%%


\subsection{Suma y diferencia de vectores}

\subsubsection*{Suma de vectores}

La suma de dos vectores libres $\vu$ y $\vv$ es otro vector obtenido de la siguiente manera: $\vu + \vv$. El protocolo para sumar vectores es el siguiente:
\begin{enumerate}
    \item Trasladamos el vector $\vv$ a continuación de $\vu$, haciendo coincidir el origen de $\vv$ con el extremo de $\vu$.

    \item El origen de la suma $\vu + \vv$ es el origen de $\vu$.

    \item El extremo de la suma $\vu + \vv$ es el extremo de $\vv$.
\end{enumerate}
Es decir, el vector $\vu + \vv$ va desde el origen de $\vu$ hasta el extremo de $\vv$. Gráficamente esto se ve así:
\begin{figure}[H]
    \centering
    \includegraphics[width=0.4\linewidth]{Imagenes/U4/fig12.png}
    \caption{}
    \label{fig:4.12}
\end{figure}
Las propiedades que cumplen la suma de vectores la veremos más adelante, pero adelantaremos dos para construir la suma de dos o más vectores:
\begin{enumerate}
    \item La suma de vectores es \Emph{asociativa}: $\forall\ \vec{a}, \vec{b}, \vec{c} \in\VV: \vec{a} + (\vec{b} + \vec{c}) = (\vec{a} + \vec{b}) + \vec{c}$

    \item La suma de vectores es \Emph{conmutativa}: $\forall\ \vec{a}, \vec{b} \in\VV: \vec{a} + \vec{b} = \vec{b} + \vec{a}$
\end{enumerate}

Estas propiedades facilitan la realización de la suma de dos o más vectores. La propiedad conmutativa permite la utilización de la regla del paralelogramo y la asociativa la posibilidad de sumar de a dos cuando tenemos dos o más vectores.

Además, si sumamos un vector libre $\vv$ con su opuesto $-\vv$, obtenemos un vector reducido a un punto (su origen y extremo coinciden), se trata del \Emph{vector nulo}. Simbólicamente: $\vv + (-\vv) = \vec{0}$.
\begin{figure}[H]
    \centering
    \includegraphics[width=0.17\linewidth]{Imagenes/U4/fig13.png}
    \caption{}
    \label{fig:4.13}
\end{figure}

%%%%%%%%%%%%%%%%%%%%%%%%%%%%%%%%%%%%%%%%%%%%%%%%%%%%%%%%%%%%

\subsubsection*{Regla del paralelogramo}

Se usa cuando necesitamos sumar dos vectores.
\begin{enumerate}
    \item Dibujamos los dos vectores $\vu$ y $\vv$ con el mismo origen.

    \item Completamos un paralelogramo trazando:
        \begin{itemize}
            \item Por el extremo del vector $\vu$, un segmento de recta paralelo al vector $\vv$.

            \item Por el extremo del vector $\vv$, un segmento de recta paralelo al vector $\vu$.
        \end{itemize}

    \item La suma de los dos vectores es la diagonal orientada del paralelogramo obtenido, que tiene su origen en el origen compartido de los vectores $\vu$ y $\vv$ (\figura{4.14} y \figura{4.15}).
\end{enumerate}

\noindent
\begin{minipage}[b]{0.45\textwidth}
    \begin{figure}[H]
        \centering
        \includegraphics[width=0.5\linewidth]{Imagenes/U4/fig14.png}
        \caption{}
        \label{fig:4.14}
    \end{figure}
\end{minipage}
\hfill
\begin{minipage}[b]{0.45\textwidth}
    \begin{figure}[H]
        \centering
        \includegraphics[width=0.5\linewidth]{Imagenes/U4/fig15.png}
        \caption{}
        \label{fig:4.15}
    \end{figure}
\end{minipage}

%%%%%%%%%%%%%%%%%%%%%%%%%%%%%%%%%%%%%%%%%%%%%%%%%%%%%%%%%%%%

\subsubsection*{Regla de la poligonal}

Conviene utilizarla cuando se suman más de dos vectores:
\begin{enumerate}
    \item A partir de un punto, se traza un vector paralelo a uno de los dados, por ejemplo, paralelo al vector $\vu$, con igual módulo y sentido que él, al que llamamos $\vu'$.
    
    \item Por el extremo de $\vu'$, trazamos un vector paralelo a $\vec{x}$, $\vec{x}'$ también con igual módulo y sentido.
    
    \item Se trabaja de igual manera con todos los otros vectores, trazando vectores equipolentes a los dados, uno a continuación del otro.
    
    \item El vector suma se obtiene uniendo el origen del primer vector con el extremo del último vector.
\end{enumerate}

\begin{figure}[H]
    \centering
    \includegraphics[width=0.4\linewidth]{Imagenes/U4/fig16.png}
    \caption{}
    \label{fig:4.16}
\end{figure}


%%%%%%%%%%%%%%%%%%%%%%%%%%%%%%%%%%%%%%%%%%%%%%%%%%%%%%%%%%%%
%%%%%%%%%%%%%%%%%%%%%%%%%%%%%%%%%%%%%%%%%%%%%%%%%%%%%%%%%%%%
%%%%%%%%%%%%%%%%%%%%%%%%%%%%%%%%%%%%%%%%%%%%%%%%%%%%%%%%%%%%


\subsubsection*{Diferencia de vectores}

La resta o diferencia entre dos vectores $\vu$ y $\vv$ se expresa como $\vu - \vv$ y se define como la suma del primero de ellos con el opuesto del segundo, es decir: $\vu - \vv = \vu + (-\vv)$.

Para representar gráficamente la diferencia $\vu - \vv$ podemos colocar a $-\vv$ a continuación de $\vu$ y unir el origen de $\vu$ con el extremo de $-\vv$, como lo indica la \figura{4.17}:
\begin{figure}[H]
    \centering
    \includegraphics[width=0.4\linewidth]{Imagenes/U4/fig17.png}
    \caption{}
    \label{fig:4.17}
\end{figure}
Usando la regla del paralelogramo, también se puede restar vectores:
\begin{figure}[H]
    \centering
    \includegraphics[width=0.2\linewidth]{Imagenes/U4/fig18.png}
    \caption{}
    \label{fig:4.18}
\end{figure}

\begin{tcolorbox}[colframe= black, boxrule=0.8pt, sharp corners]
    Por lo tanto, el vector \Emph{suma} de dos vectores coincide con una de las diagonales del paralelogramo, la otra diagonal representa la \Emph{resta} de dichos vectores.
\end{tcolorbox}

\begin{figure}[H]
    \centering
    \includegraphics[width=0.5\linewidth]{Imagenes/U4/fig19.png}
    \caption{}
    \label{fig:4.19}
\end{figure}

Para efectuar sumas o restas de tres o más vectores, el proceso es idéntico. Basta con aplicar la propiedad asociativa.

Al vector que se obtiene al sumar o restar varios vectores se le denomina \Emph{resultante}.


%%%%%%%%%%%%%%%%%%%%%%%%%%%%%%%%%%%%%%%%%%%%%%%%%%%%%%%%%%%%
%%%%%%%%%%%%%%%%%%%%%%%%%%%%%%%%%%%%%%%%%%%%%%%%%%%%%%%%%%%%
%%%%%%%%%%%%%%%%%%%%%%%%%%%%%%%%%%%%%%%%%%%%%%%%%%%%%%%%%%%%


\subsubsection*{Expresión analítica de la suma y resta de vectores}

Dados los vectores $\vu = (x_1, y_1), \vv = (x_2, y_2) \in\RR^2$, se definen la suma y la resta entre los vectores $\vu$ y $\vv$ a los vectores:
    \begin{align*}
        \vu + \vv &= (x_1, y_1) + (x_2, y_2) = (x_1 + x_2,\ y_1 + y_2) \\
        \vu - \vv &= (x_1, y_1) - (x_2, y_2) = (x_1 - x_2,\ y_1 - y_2)
    \end{align*}
De manera similar, se definen estas operaciones en el espacio tridimensional. Sean $\vu = (x_1, y_1, z_1), \vv = (x_2, y_2, z_2) \in\RR^3$:
    \begin{align*}
        \vu + \vv &= (x_1, y_1, z_1) + (x_2, y_2, z_2) = (x_1 + x_2,\ y_1 + y_2,\ z_1 + z_2)\\
        \vu - \vv &= (x_1, y_1, z_1) - (x_2, y_2, z_2) = (x_1 - x_2,\ y_1 - y_2,\ z_1 - z_2)
    \end{align*}
\begin{tcolorbox}[colframe= black, boxrule=0.8pt, sharp corners]
    La suma (resta) de dos vectores es otro vector cuyas componentes es la suma (resta) de sus componentes homólogas.
\end{tcolorbox}

\begin{ejemplo}
    Sean los vectores $\vu = (-2,3)$ y $\vv = (1,2)$, entonces:
        \begin{align*}
            \vu +\vv &= (-2,3) + (1,2) = (-2+1,\ 3+2) = (-1, 5) \\
            \vu -\vv &= (-2,3) - (1,2) = (-2-1,\ 3-2) = (-3, 1)
        \end{align*}
    Gráficamente, esto es:
    
    \noindent
    \begin{minipage}[b]{0.45\textwidth}
        \begin{figure}[H]
            \centering
            \includegraphics[width=0.4\linewidth]{Imagenes/U4/fig20.png}
        \end{figure}
    \end{minipage}
    \hfill
    \begin{minipage}[b]{0.45\textwidth}
        \begin{figure}[H]
            \centering
            \includegraphics[width=0.5\linewidth]{Imagenes/U4/fig21.png}
        \end{figure}
    \end{minipage}
\end{ejemplo}

%%%%%%%%%%%%%%%%%%%%%%%%%%%%%%%%%%%%%%%%%%%%%%%%%%%%%%%%%%%%
%%%%%%%%%%%%%%%%%%%%%%%%%%%%%%%%%%%%%%%%%%%%%%%%%%%%%%%%%%%%
%%%%%%%%%%%%%%%%%%%%%%%%%%%%%%%%%%%%%%%%%%%%%%%%%%%%%%%%%%%%


\subsubsection*{Propiedades}

Sean los vectores $\vec{v_1}, \vec{v_2}, \vec{v_3} \in\VV$, siendo $\VV$ el conjunto de todos los vectores libres, y sea $+$ una operación que llamaremos \Emph{suma algebraica de vectores}, entonces, se verifica que:
\begin{enumerate}
    \item \Emph{Ley de composición interna:} $\forall\ \vec{v_1}, \vec{v_2} \in\VV: \vec{v_1} + \vec{v_2} \in\VV$
    \begin{figure}[H]
        \centering
        \includegraphics[width=0.2\linewidth]{Imagenes/U4/fig22.png}
        \caption{}
        \label{fig:4.22}
    \end{figure}

    \item \Emph{Propiedad asociativa:} $\forall\ \vec{v_1}, \vec{v_2}, \vec{v_3} \in\VV: \vec{v_1} + (\vec{v_2} + \vec{v_3}) = (\vec{v_1} + \vec{v_2}) + \vec{v_3}$
    \begin{figure}[H]
        \centering
        \includegraphics[width=0.3\linewidth]{Imagenes/U4/fig23.png}
        \caption{}
        \label{fig:4.23}
    \end{figure}

    \item \Emph{Existencia del elemento neutro:} $\exists\ \vec{0}\in\VV\ /\ \forall\ \vv\in\VV: \vv + \vec{0} = \vv$
    \begin{figure}[H]
        \centering
        \includegraphics[width=0.15\linewidth]{Imagenes/U4/fig24.png}
        \caption{}
        \label{fig:4.24}
    \end{figure}

    \item \Emph{Existencia del elemento simétrico u opuesto:} $\forall\ \vv\in\VV: \exists\ (-\vv)\in\VV\ /\ \vv + (-\vv) = \vec{0}$
    \begin{figure}[H]
        \centering
        \includegraphics[width=0.15\linewidth]{Imagenes/U4/fig25.png}
        \caption{}
        \label{fig:4.25}
    \end{figure}
    
    \item \Emph{Propiedad conmutativa}: $\forall\ \vec{v_1}, \vec{v_2} \in\VV: \vec{v_1} + \vec{v_2} = \vec{v_2} + \vec{v_1}$
    \begin{figure}[H]
        \centering
        \includegraphics[width=0.4\linewidth]{Imagenes/U4/fig26.png}
        \caption{}
        \label{fig:4.26}
    \end{figure}
\end{enumerate}


%%%%%%%%%%%%%%%%%%%%%%%%%%%%%%%%%%%%%%%%%%%%%%%%%%%%%%%%%%%%
%%%%%%%%%%%%%%%%%%%%%%%%%%%%%%%%%%%%%%%%%%%%%%%%%%%%%%%%%%%%
%%%%%%%%%%%%%%%%%%%%%%%%%%%%%%%%%%%%%%%%%%%%%%%%%%%%%%%%%%%%


\subsection{Producto de un escalar por un vector}

\noindent
\begin{minipage}[t]{0.6\textwidth}
    Se define el producto de un número real $k$ por un vector $\vec{u}$ como el vector que tiene:
    \begin{enumerate}
        \item \Emph{Dirección:} la misma que $\vec{u}$
        
        \item \Emph{Sentido:} el mismo que $\vec{u}$, si $k$ es positivo y opuesto al de $\vec{u}$, si $k$ es negativo. (\figura{4.27})
        
        \item \Emph{Módulo:} el módulo es el valor absoluto de $k$ multiplicado por el módulo de $\vec{u}$. Es decir:
            \begin{equation*}
                |k \vec{u}| = |k| . |\vec{u}|
            \end{equation*}
    \end{enumerate}
\end{minipage}
\hfill
\begin{minipage}[t]{0.35\textwidth}
    \vspace{-.7cm}
    \begin{figure}[H]
        \centering
        \includegraphics[width=0.8\linewidth]{Imagenes/U4/fig27.png}
        \caption{}
        \label{fig:4.27}
    \end{figure}
\end{minipage}
\vspace{0.2cm}

Sean $\vv\in\VV$, un vector libre y $k\in\RR$ un escalar, se distinguen los siguientes casos del producto de un escalar por un vector:
\begin{enumerate}[label=(\roman*)]
    \item $k>1:$ En este caso, el producto $k\cdot\vv$ es un vector con el mismo sentido y dirección que $\vv$ pero de módulo mayor que $\vv$.
    \begin{figure}[H]
        \centering
        \includegraphics[width=0.25\linewidth]{Imagenes/U4/fig28.png}
        \caption{}
        \label{fig:4.28}
    \end{figure}

    \item $k<-1:$ En este caso, el producto $k\cdot\vv$ es un vector con la misma dirección que $\vv$ pero de sentido opuesto a $\vv$ y de módulo mayor que el mismo.
    \begin{figure}[H]
        \centering
        \includegraphics[width=0.25\linewidth]{Imagenes/U4/fig29.png}
        \caption{}
        \label{fig:4.29}
    \end{figure} 

    \item $-1<k<1:$ En este caso, el producto $k\cdot\vv$ es un vector con la misma dirección que $\vv$ pero su módulo es menor que $\vv$ y su sentido dependerá del signo de $k$, como se vio en los dos casos anteriores.
    \begin{figure}[H]
        \centering
        \includegraphics[width=0.3\linewidth]{Imagenes/U4/fig30.png}
        \caption{}
        \label{fig:4.30}
    \end{figure}

    \item $k=0:$ En este caso, el producto $k\cdot\vv$ se reduce al vector nulo. Es decir: $k\cdot\vv = \vec{0}$
\end{enumerate}

En resumen, multiplicar un vector por un número real $k$ equivale a alargar (o encoger) su módulo tantas veces como indica el valor absoluto de $k$, e invertir su sentido si $k$ es negativo. El número real $k$ por el que se multiplica un vector recibe el nombre de \Emph{escalar}.

\begin{tcolorbox}[colframe= black, boxrule=0.8pt, sharp corners]
    Dos \Emph{vectores} son \Emph{paralelos} si, y solo si, son múltiplos escalares uno del otro.    
\end{tcolorbox}

Analíticamente, si $k \in \RR$ y $\vv = (x, y) \in \RR^2$, el producto escalar vector se define de la siguiente manera:
    \begin{equation*}
        k\cdot\vv = (kx,\ ky)
    \end{equation*}
Es decir, cada componente del vector $\vv$ queda multiplicada por el escalar $k$. Análogamente en $\RR^3$:
\begin{equation*}
    k\cdot\vv = (x, y, z) = (kx,\ ky,\ kz)
\end{equation*}

\begin{ejemplo}
    \begin{minipage}[t]{0.6\textwidth}
        Sean los vectores $\vv = (2,1),\ \vu = (2,2)\in\RR^2$, sean $k_1 = 2$ y $k_2 = -1$. Entonces:
            \begin{align*}
                \\
                k_1\cdot\vv &= 2\cdot(2, 1) = (4,2) \\
                \\
                k_2\cdot\vu &= (-1)\cdot(2, 2) = (-2,-2)
            \end{align*}
    \end{minipage}
    \hfill
    \begin{minipage}[t]{0.35\textwidth}
        \vspace{-1cm}
        \begin{figure}[H]
            \centering
            \includegraphics[width=.8\linewidth]{Imagenes/U4/fig31.png}
            \caption{}
            \label{fig:4.31}
        \end{figure}
    \end{minipage}
\end{ejemplo}

%%%%%%%%%%%%%%%%%%%%%%%%%%%%%%%%%%%%%%%%%%%%%%%%%%%%%%%%%%%%
%%%%%%%%%%%%%%%%%%%%%%%%%%%%%%%%%%%%%%%%%%%%%%%%%%%%%%%%%%%%
%%%%%%%%%%%%%%%%%%%%%%%%%%%%%%%%%%%%%%%%%%%%%%%%%%%%%%%%%%%%


\subsubsection*{Propiedades}

\begin{enumerate}
    \item \Emph{Ley de composición externa:} $\forall\ k \in\RR, \forall\ \vv \in\VV: k \cdot \vv \in\VV$
    
    \item \Emph{Propiedad asociativa mixta:} $\forall\ k_1, k_2 \in\RR, \forall\ \vv \in\VV: (k_1\ k_2) \cdot \vv = k_1 (k_2\cdot \vv)$
    
    \item \Emph{Propiedad distributiva del producto de un vector respecto a la suma de escalares:}\\
    $\forall\ k_1, k_2 \in \RR, \forall\ \vv \in\VV: (k_1 + k_2) \cdot \vv = k_1\cdot \vv + k_2\cdot \vv$
    
    \item \Emph{Propiedad distributiva del producto de un escalar respecto a la suma de vectores:} \\
    $\forall\ k \in \RR, \forall\ \vec{v_1}, \vec{v_2} \in \VV: k \cdot (\vec{v_1} + \vec{v_2}) = k\cdot \vec{v_1} + k\cdot \vec{v_2}$
    
    \item \Emph{El elemento neutro para el producto de un escalar por un vector:}\\ 
    $\exists\ 1 \in \RR\ /\ \forall\ \vv \in\VV: 1 \cdot \vv = \vv$
\end{enumerate}

Por las propiedades de la suma de vectores y el producto escalar vector, se tiene que la cuaterna $(\VV, +, \RR, \cdot)$ es un \Emph{espacio vectorial}. De esta manera $(\RR^2, +, \RR, \cdot)$ y $(\RR^3, +, \RR, .)$ son
espacios vectoriales reales.


%%%%%%%%%%%%%%%%%%%%%%%%%%%%%%%%%%%%%%%%%%%%%%%%%%%%%%%%%%%%
%%%%%%%%%%%%%%%%%%%%%%%%%%%%%%%%%%%%%%%%%%%%%%%%%%%%%%%%%%%%
%%%%%%%%%%%%%%%%%%%%%%%%%%%%%%%%%%%%%%%%%%%%%%%%%%%%%%%%%%%%


\subsection{Vector determinado por dos puntos}

\subsubsection* {En $\RR^2$}
Los puntos $A(x_1, y_1)$  y  $B(x_2, y_2)$, determinan los vectores $\vec{OA}, \vec{OB}$ y $\vec{AB}$. Por la suma y resta de vectores es:

\noindent
\begin{minipage}[t]{0.6\textwidth}
        \begin{equation*}
            \vec{OA} + \vec{AB} = \vec{OB} \iff \vec{AB} = \vec{OB} - \vec{OA}
        \end{equation*}    
    Reemplazando por las componentes de los vectores:
        \begin{equation*}
            \vec{AB} = (x_2, y_2) - (x_1, y_1)
        \end{equation*}
    Por definición de resta de vectores,
    tenemos:
        \begin{equation*}
            \vec{AB} = (x_2 - x_1,\ y_2 - y_1)
        \end{equation*}
\end{minipage}
\hfill
\begin{minipage}[t]{0.35\textwidth}
    \vspace{-1cm}
    \begin{figure}[H]
        \centering
        \includegraphics[width=0.8\linewidth]{Imagenes/U4/fig32.png}
        \caption{}
        \label{fig:4.32}
    \end{figure}
\end{minipage}
\vspace{0.2cm}

\noindent El módulo de este vector será: $|\vec{AB}| = \sqrt{(x_2 - x_1)^2 + (y_2 - y_1)^2}$.

\begin{tcolorbox}[colframe= black, boxrule=0.8pt, sharp corners]
    Esta fórmula también define la distancia entre dos puntos del plano, es decir, distancia entre $A$ y $B$:
        \begin{equation*}
            d(A,B) = \sqrt{(x_2 - x_1)^2 + (y_2 - y_1)^2}
        \end{equation*}
\end{tcolorbox}

%%%%%%%%%%%%%%%%%%%%%%%%%%%%%%%%%%%%%%%%%%%%%%%%%%%%%%%%%%%%

\subsubsection*{En $\RR^3$}

Si trabajamos de igual forma que en el plano, tendremos que las componentes de un vector determinado por dos puntos son:
    \begin{equation*}
        \vec{AB} = (x_2 - x_1,\ y_2 - y_1,\ z_2 - z_1)
    \end{equation*}
Siendo el módulo de dicho vector: $|\vec{AB}| = \sqrt{(x_2 - x_1)^2 + (y_2 - y_1)^2 + (z_2 - z_1)^2}$.

\begin{tcolorbox}[colframe= black, boxrule=0.8pt, sharp corners]
    Al calcular el módulo del vector determinado por $\vec{A}$ y $\vec{B}$, estamos calculando la longitud del segmento $\overline{AB}$ y la distancia entre dos puntos en el espacio $A$ y $B$:
        \begin{equation*}
            d(A, B) = \sqrt{(x_2 - x_1)^2 + (y_2 - y_1)^2 + (z_2 - z_1)^2}
        \end{equation*}
\end{tcolorbox}

\noindent
\begin{minipage}[b]{0.45\textwidth}
    \vspace{-.3cm}
    \begin{figure}[H]
        \centering
        \includegraphics[width=.9\linewidth]{Imagenes/U4/fig33.png}
        \caption{}
        \label{fig:4.33}
    \end{figure}
\end{minipage}
\hfill
\begin{minipage}[b]{0.45\textwidth}
    \vspace{-.4cm}
    \begin{figure}[H]
        \centering
        \includegraphics[width=.8\linewidth]{Imagenes/U4/fig34.png}
        \caption{}
        \label{fig:4.34}
    \end{figure}
\end{minipage}


%%%%%%%%%%%%%%%%%%%%%%%%%%%%%%%%%%%%%%%%%%%%%%%%%%%%%%%%%%%%
%%%%%%%%%%%%%%%%%%%%%%%%%%%%%%%%%%%%%%%%%%%%%%%%%%%%%%%%%%%%
%%%%%%%%%%%%%%%%%%%%%%%%%%%%%%%%%%%%%%%%%%%%%%%%%%%%%%%%%%%%


\subsection{Coordenadas del punto medio de un segmento}

Consideramos dos puntos sobre un vector $\vec{P_1 P_2}$, siendo $P_1(x_1, y_1)$ y $P_2(x_2, y_2)$. El punto medio del segmento $\vec{P_1P_2}$ es $P_m(x_m, y_m)$, por lo tanto, se deduce que: $\overline{P_1 P_m} = \overline{P_m P_2}$.

\noindent
\begin{minipage}[t]{0.35\textwidth}
    \vspace{-.3cm}
    \begin{figure}[H]
        \centering
        \includegraphics[width=0.8\linewidth]{Imagenes/U4/fig35.png}
        \caption{}
        \label{fig:4.35}
    \end{figure}
\end{minipage}
\hfill
\begin{minipage}[t]{0.6\textwidth}
        \begin{align*}
            &\implies \begin{cases}
                x_m - x_1 = x_2 - x_m \\
                y_m - y_1 = y_2 - y_m
            \end{cases} \\
            &\implies \begin{cases}
                2x_m = x_2 + x_1 \iff x_m = (x_2 + x_1)/2\\
                2y_m =\, y_2 + y_1 \iff y_m =\ (y_2 + y_1)/2
            \end{cases}
        \end{align*}
    Y por lo tanto:
        \begin{equation*}
            P_m = (x_m, y_m) = \left(\frac{x_1 + x_2}{2},\ \frac{y_1 + y_2}{2}\right)
        \end{equation*}
\end{minipage}

\begin{ejemplo}
    \noindent
    \begin{minipage}[t]{0.6\textwidth}
        Determinar el punto medio del segmento de extremos $P_1(1, 2)$ y $P_2(7, 6)$.
        \vspace{.5cm}
            \begin{equation*}
                \begin{rcases}
                    x_m &= \dfrac{1+7}{2} = 4 \\
                    \\
                    y_m &= \dfrac{2+6}{2} = 4
                \end{rcases}
                P(4, 4)
            \end{equation*}
    \end{minipage}
    \hfill
    \begin{minipage}[t]{0.35\textwidth}
        \vspace{-1cm}
        \begin{figure}[H]
            \centering
            \includegraphics[width=.9\linewidth]{Imagenes/U4/fig36.png}
            \caption{}
            \label{fig:4.36}
        \end{figure}
    \end{minipage}
    \vspace{0.2cm}
\end{ejemplo}

En el espacio tridimensional, las coordenadas del \Emph{punto medio del segmento} que une los puntos $P_1(x_1, y_1, z_1)$ y $P_2(x_2, y_2, z_2)$ son:
    \begin{equation*}
        P_m = (x_m, y_m, z_m) = \left(\frac{x_1 + x_2}{2},\ \frac{y_1 + y_2}{2},\ \frac{z_1 + z_2}{2}\right)
    \end{equation*}

\begin{preguntabox}[Resolver:]
    Dados los puntos $A$ y $B$, hallar las coordenadas del punto $C$ de manera que $B$ sea el punto medio del segmento $\overline{AC}$:
    \begin{enumerate}
        \item $A(1,1)$ y $B(2; 2,5)$
        \item $A(1, -2)$ y $B(3, -3)$
    \end{enumerate}
\end{preguntabox}


%%%%%%%%%%%%%%%%%%%%%%%%%%%%%%%%%%%%%%%%%%%%%%%%%%%%%%%%%%%%
%%%%%%%%%%%%%%%%%%%%%%%%%%%%%%%%%%%%%%%%%%%%%%%%%%%%%%%%%%%%
%%%%%%%%%%%%%%%%%%%%%%%%%%%%%%%%%%%%%%%%%%%%%%%%%%%%%%%%%%%%


\subsection{Combinación lineal}

Dado un conjunto $A = \{\vec{v_1}, \vec{v_2}, \dots, \vec{v_n}\} \subset \VV$ y sean los escalares $k_1, k_2, \dots, k_n$, se llama \Emph{combinación lineal} (C.L.), a todo vector de la forma:
    \begin{equation*}
        \vv = k_1 \vec{v_1} + k_2 \vec{v_2} + k_3 \vec{v_3} + \dots + k_n \vec{v_n}
    \end{equation*}
Utilizando el símbolo de sumatoria, podemos escribir esta suma de la siguiente manera:
\begin{equation*}
    \vv = \sum_{i=1}^{n} k_i \vec{v_i} \quad\text{ con } k_i \in \RR \land \vec{v_i} \in A,\ A \subset \VV
\end{equation*}
Una combinación lineal de dos o más vectores es el vector que se obtiene al sumar los productos entre los vectores y escalares dados. Los escalares se denominan \Emph{coeficientes} de la combinación lineal.

\begin{tcolorbox}[colframe= black, boxrule=0.8pt, sharp corners]
    En otras palabras, diremos que un vector $\vv$ es combinación lineal de los vectores del conjunto $A$ si, y solo si, existen escalares $k_1, k_2, \dots, k_n$, tales que $\vv$ puede expresarse como
        \begin{equation*}
            \vv = \sum_{i=1}^{n} k_i \vec{v_i}
        \end{equation*}
\end{tcolorbox}

\begin{ejemplo}
    Dado el conjunto $A = \{ \vv, \vu \}$, con $\vv = (1,2), \vu = (2,1)$, y los escalares $k_1 = 3,\ k_2 = 2$, podemos expresar al vector $\vec{w} = k_1\cdot\vv + k_2\cdot\vu = 3(1,2) + 2(2,1) = (7,8)$.
    \begin{figure}[H]
        \centering
        \includegraphics[width=0.3\linewidth]{Imagenes/U4/fig37.png}
        \caption{}
        \label{fig:4.37}
    \end{figure}
\end{ejemplo}


%%%%%%%%%%%%%%%%%%%%%%%%%%%%%%%%%%%%%%%%%%%%%%%%%%%%%%%%%%%%
%%%%%%%%%%%%%%%%%%%%%%%%%%%%%%%%%%%%%%%%%%%%%%%%%%%%%%%%%%%%
%%%%%%%%%%%%%%%%%%%%%%%%%%%%%%%%%%%%%%%%%%%%%%%%%%%%%%%%%%%%


\subsection{Dependencia e independencia lineal}

\subsubsection*{Dependencia lineal}

\begin{definitionbox}{Vectores Linealmente Dependientes}{4.2}
    Un conjunto de vectores $A = \{\vv_1, \vv_2, \dots, \vv_n\} \subset \VV$ es \Emph{linealmente dependientes} (L.D.) si, y solo si, existe una combinación lineal de ellos que dé como resultado el vector nulo, sin que sean cero todos los coeficientes de la combinación lineal.
        \begin{equation*}
            A \text{ es L.D.} \iff \exists\ i\ /\ k_1 \vec{v_1} + k_2 \vec{v_2} + \dots + k_n \vec{v_n} = \sum_{i=1}^{n} k_i \vec{v_i} = \vec{0}\ \land\ k_i\neq0
        \end{equation*}
\end{definitionbox}


%%%%%%%%%%%%%%%%%%%%%%%%%%%%%%%%%%%%%%%%%%%%%%%%%%%%%%%%%%%%


\subsubsection*{Propiedades de vectores L.D.}

\begin{enumerate}
    \item Si un conjunto finito y no vacío de vectores son linealmente dependientes, entonces al menos uno de ellos se puede expresar como combinación lineal de los demás.
    
    \begin{demostracionbox}        
        Sea $A = \{\vec{v_1}, \vec{v_2}, \dots, \vec{v_n}\} \subset \VV$ un conjunto de vectores L.D. de $\VV$. Por definición de dependencia lineal, existen escalares $k_1, k_2, \dots, k_n$ no todos nulos tales que:        
            \begin{equation*}
                k_1 \vec{v_1} + k_2 \vec{v_2} + k_3 \vec{v_3} + \dots + k_n \vec{v_n} = \vec{0}
            \end{equation*}
        Como al menos uno de los coeficientes debe ser distinto de cero, podemos suponer, sin perder generalidad, que $k_1 \neq 0$, por lo que existe su inverso. Despejando $\vec{v_1}$:
        
        \begin{equation*}
            \vec{v_1} = -\frac{k_2}{k_1}\ \vec{v_2} - \frac{k_3}{k_1}\ \vec{v_3} - \dots - \frac{k_n}{k_1}\ \vec{v_n}
        \end{equation*}
    \end{demostracionbox}


    \item Si un vector de $A = \{\vec{v_1}, \vec{v_2}, \dots, \vec{v_n}\}$ es C.L. de los demás vectores de $A$, entonces los vectores de $A$ son L.D.
    
    \begin{demostracionbox}        
        Suponemos que $\vec{v_1}$ es C.L. de los vectores de $A$:
            \begin{equation*}
                k_2 \vec{v_2} + k_3 \vec{v_3} + \dots + k_n \vec{v_n} = \vec{v_1}
            \end{equation*}
        Sumando a ambos miembros el opuesto del vector $\vec{v_1}$:
            \begin{equation*}
                -\vec{v_1} + k_2 \vec{v_2} + k_3 \vec{v_3} + \dots + k_n \vec{v_n} = \vec{0}
            \end{equation*}        
        Teniendo en cuenta que:
            \begin{equation*}
                -\vec{v_1} = (-1)\cdot\vec{v_1} \implies (-1)\cdot\vec{v_1} + k_2 \vec{v_2} + k_3 \vec{v_3} + \dots + k_n \vec{v_n} = \vec{0}
            \end{equation*}
        En la última igualdad, se observa que al menos existe un coeficiente distinto de cero, es decir $-1$, en consecuencia, $A$ es L.D.
    \end{demostracionbox}

    \begin{ejemplo}
        Si al vector $\vv = 2\cdot \vec{v_1} + 3\cdot \vec{v_2} + 4\cdot \vec{v_3}$, le sumamos $(-\vec{v})$, obtenemos:
            \begin{equation*}
                2\cdot \vec{v_1} + 3\cdot \vec{v_2} + 4\cdot \vec{v_3} + (-1)\vec{v} = \vec{0}
            \end{equation*}
        \begin{figure}[H]
            \centering
            \includegraphics[width=0.35\textwidth]{Imagenes/U4/fig38.png}
            \caption{}
            \label{fig:4.38}
        \end{figure}        
        Entonces, el conjunto $\{\vec{v}, \vec{v}_1, \vec{v}_2, \vec{v}_3\}$ es un conjunto de vectores linealmente dependiente.
    \end{ejemplo}


    \item Dos vectores del plano son linealmente dependientes sí, y sólo sí, son paralelos.


    \item Dos vectores libres del plano $\vu = (x_1, y_1)$ y $\vv = (x_2, y_2)$ son linealmente dependientes si sus componentes son proporcionales.    
        \begin{equation*}
            \vu = k\cdot \vv \implies (x_1, y_1) = k\cdot(x_2, y_2) \implies \frac{x_1}{x_2} = \frac{y_1}{y_2} = k
        \end{equation*}
\end{enumerate}


%%%%%%%%%%%%%%%%%%%%%%%%%%%%%%%%%%%%%%%%%%%%%%%%%%%%%%%%%%%%


\subsubsection*{Independencia lineal}

\begin{definitionbox}{Vectores Linealmente Independientes}{4.3}
    El conjunto de vectores $A = \{\vec{v_1}, \vec{v_2}, \dots, \vec{v_n}\} \subset \VV$, es \Emph{linealmente independientes} (L.I.) si, y solo si, la única manera de obtener el vector nulo como combinación lineal de ellos, es que \Emph{todos} los coeficientes sean ceros.
        \begin{equation*}
            A \text{ es L.I.} \iff k_1 \vec{v_1} + k_2 \vec{v_2} + \dots + k_n \vec{v_n} = \sum_{i=1}^{n} k_i \vec{v_i} = \vec{0}\ \land\  k_i = 0,\ \forall\ i 
        \end{equation*}
\end{definitionbox}

%%%%%%%%%%%%%%%%%%%%%%%%%%%%%%%%%%%%%%%%%%%%%%%%%%%%%%%%%%%%


\subsubsection*{Propiedades de vectores L.I.}

\begin{enumerate}
    \item Un conjunto de vectores es L.I. si, y solo si, ningún vector se puede expresar como combinación lineal de los demás.
    
    \item Si un vector es combinación lineal de un conjunto de vectores linealmente independiente, entonces dicha combinación lineal es única.
    
    \item Dos vectores del plano son L.I. si, y solo si, no son paralelos.
    
    \item Dos vectores del plano son L.I. si, y solo si, sus componentes no son proporcionales.
\end{enumerate}
Para investigar la dependencia o independencia lineal de un conjunto de vectores, se propone una combinación lineal de dicho conjunto, con escalares a determinar, que sea igual al vector nulo. Si algún escalar es no nulo, el conjunto es L.D., mientras que, si todos los escalares son necesariamente nulos, el conjunto es L.I.


%%%%%%%%%%%%%%%%%%%%%%%%%%%%%%%%%%%%%%%%%%%%%%%%%%%%%%%%%%%%


\subsubsection*{En $\RR^2$:}

\begin{itemize}
    \item Dos vectores paralelos son L.D.
        \begin{figure}[H]
            \centering
            \includegraphics[width=0.3\linewidth]{Imagenes/U4/fig39.png}
            %\caption{}
            \label{fig:4.39}
        \end{figure}

    \item Dos vectores no paralelos son L.I.
        \begin{figure}[H]
            \centering
            \includegraphics[width=0.35\linewidth]{Imagenes/U4/fig40.png}
            %\caption{}
            \label{fig:4.40}
        \end{figure}

    \item Tres vectores son L.D.
        \begin{figure}[H]
            \centering
            \includegraphics[width=0.4\linewidth]{Imagenes/U4/fig41.png}
            %\caption{}
            \label{fig:4.41}
        \end{figure}
\end{itemize}



\subsubsection*{En $\RR^3$:}

\begin{itemize}
    \item Tres vectores no coplanares son L.I: $k_1 \vec{v_1} + k_2 \vec{v_2} + k_3 \vec{v_3} = \vec{0} \iff k_1 = k_2 = k_3 = 0$
        \begin{figure}[H]
            \centering
            \includegraphics[width=0.3\linewidth]{Imagenes/U4/fig42.png}
            %\caption{}
            \label{fig:4.42}
        \end{figure}

    \item Cuatro vectores son L.D:
        \begin{equation*}
            \begin{aligned}
                \vec{v_4} = \vec{v_1} + &2\cdot \vec{v_2} + \vec{v_3} \\
                \vec{v_1} + 2\cdot \vec{v_2} + \vec{v_3} &+ (-1)\cdot \vec{v_4} = \vec{0}
            \end{aligned}
        \end{equation*}
        \begin{figure}[H]
            \centering
            \includegraphics[width=0.3\linewidth]{Imagenes/U4/fig43.png}
            %\caption{}
            \label{fig:4.43}
        \end{figure}
\end{itemize}



\begin{preguntabox}[Resolver:]
    \begin{enumerate}
        \item Si $\vu$ y $\vv$ son dos vectores que verifican $3\vu + \frac{1}{3}\vv = \vec{0}$ entonces, ¿$\vu$ y $\vv$ son linealmente dependientes o independientes? Justificar.

        \item Dados los vectores $\vv = (4, 6, m)$ y $\vec{w} = (2, k, 1)$ determinar, si existen, los valores de $m$ y $k$ para que los vectores $\vv$ y $\vec{w}$ sean linealmente dependientes.
    \end{enumerate}
\end{preguntabox}

\subsection{Guía de repaso y desarrollo teórico}

A continuación, encontrarás una serie de preguntas y consignas que te ayudarán a repasar los temas principales de la unidad. Algunas te invitan a recordar definiciones o propiedades, y otras a desarrollar los conceptos en forma completa, clara y ordenada.

\subsubsection*{Preguntas}

\begin{enumerate}
    \item Definir vector geométrico. Mencionar sus características y explicar brevemente cada una de ellas. Realizar un gráfico señalando los elementos que lo componen.
    
    \item Deducir la expresión del módulo o norma de un vector en el plano, en función de sus componentes, a partir del teorema de Pitágoras. Proponer un ejemplo.
    
    \item Escribir la condición analítica de igualdad entre dos vectores del espacio, expresando la relación entre sus componentes.
    
    \item Si dos vectores son iguales, ¿pueden tener distintos puntos de aplicación? Justificar la respuesta desde la interpretación geométrica del concepto.
    
    \item Si la suma de dos vectores es el vector nulo, ¿qué relación hay entre ellos? Proponer un ejemplo.
    
    \item Dados dos vectores del plano, escribir la expresión analítica de su suma y representarla gráficamente utilizando la regla del paralelogramo. Explicar brevemente el procedimiento realizado.
    
    \item Si al multiplicar un vector no nulo por un número real $k$ se obtiene otro vector con igual dirección y sentido opuesto, ¿qué puede afirmarse sobre el valor de $k$?
    
    \item Deducir la fórmula del vector determinado por dos puntos y, a partir de ella, la expresión de la distancia entre los puntos en el plano. Representar gráficamente la situación, identificando los elementos que intervienen en la deducción.
    
    \item Sean los puntos $P_1(1,a)$ y $P_2(b,6)$. Hallar las coordenadas $a$ y $b$, sabiendo que el punto medio del segmento $\overline{P_1P_2}$ es $P_m(2,1)$.
    
    \item ¿Cuándo es posible expresar a un vector como combinación lineal de otros?
    
    \item Escribir la definición coloquial y simbólica de combinación lineal de vectores. Representar gráficamente una combinación lineal de dos vectores en el plano.
    
    \item Sea el conjunto $A=\{\vec{u},\vec{v},\vec{w}\}$, siendo $\vec{u}=\alpha\vec{v}+\beta\vec{w}$. Determinar si el conjunto $A$ es linealmente dependiente o independiente. Justificar la elección de la respuesta.
\end{enumerate}
